\subsection{Morita Equivalence of Algebras}

PROPOSITION 13. --- a) If two k-algebras are Morita equivalent, then their centers are isomorphic k-algebras.

\begin{enumerate}
\def\labelenumi{\alph{enumi})}
\setcounter{enumi}{1}
\item
  Two commutative \(k\) -algebras are Morita equivalent if and only if they are isomorphic.
\item
  Two k-algebras that are fields are Morita equivalent if and only if they are isomorphic.
\item
  For \(i = 1,2\) , let \(\mathrm{A}_i\) and \(\mathrm{B}_i\) be Morita equivalent \(k\) -algebras and \(\mathrm{P}_i\) an invertible \((\mathrm{A}_i,\mathrm{B}_i)_k\) -bimodule. Set \(\mathrm{A} = \mathrm{A}_1\otimes_k\mathrm{A}_2\) , \(\mathrm{B} = \mathrm{B}_1\otimes_k\mathrm{B}_2\) , and \(\mathrm{P} = \mathrm{P}_1 \otimes_k \mathrm{P}_2\) . The \(k\) -algebras A and B are Morita equivalent, and P is an invertible \((\mathrm{A}, \mathrm{B})_k\) -bimodule.
\item
  If A and B are Morita equivalent \(k\) -algebras and \(k'\) is a commutative \(k\) -algebra, then the \(k'\) -algebras \(\mathrm{A}_{(k')}\) and \(\mathrm{B}_{(k')}\) are Morita equivalent.
\end{enumerate}

Assertion a) follows from Corollary 1 of VIII, p.~102, and implies b).

Let \(\mathrm{K}\) and \(\mathrm{L}\) be \(k\) -algebras that are fields, and let \(\mathrm{P}\) be an invertible \((\mathrm{K},\mathrm{L})_k\) -bimodule. The right \(\mathrm{L}\) -vector space \(\mathrm{P}\) is a simple module (VIII, p.~104), hence of dimension 1, so the \(k\) -algebras \(\operatorname{End}_{\mathrm{L}}(\mathrm{P})\) and \(\mathrm{L}\) are isomorphic. By VIII, p.~101, Theorem 1, the mapping \(a\mapsto a_{\mathrm{P}}\) from \(\mathrm{K}\) to \(\operatorname{End}_{\mathrm{L}}(\mathrm{P})\) is an isomorphism. Therefore, the fields \(\mathrm{K}\) and \(\mathrm{L}\) are isomorphic over \(k\) ; assertion c) follows.

Under the assumptions of d), let \(Q_{i}\) (i = 1, 2) be a \((\mathrm{B}_{i}, \mathrm{A}_{i})_{k}\) -bimodule inverse to \(P_{i}\) . Denote the \((\mathrm{B}, \mathrm{A})_{k}\) -bimodule \(Q_{1} \otimes_{k} Q_{2}\) by Q. Consider the canonical k-linear isomorphism \((\mathrm{P}_{1} \otimes_{k} \mathrm{P}_{2}) \otimes_{k} (\mathrm{Q}_{1} \otimes_{k} \mathrm{Q}_{2}) \to (\mathrm{P}_{1} \otimes_{k} \mathrm{Q}_{1}) \otimes_{k} (\mathrm{P}_{2} \otimes_{k} \mathrm{Q}_{2})\) ; when passing to the quotient, it defines a morphism

\[
\left(\mathrm{P} _ {1} \otimes_ {k} \mathrm{P} _ {2}\right) \otimes_ {\mathrm{B}} \left(\mathrm{Q} _ {1} \otimes_ {k} \mathrm{Q} _ {2}\right)\rightarrow \left(\mathrm{P} _ {1} \otimes_ {\mathrm{B} _ {1}} \mathrm{Q} _ {1}\right) \otimes_ {k} \left(\mathrm{P} _ {2} \otimes_ {\mathrm{B} _ {2}} \mathrm{Q} _ {2}\right)
\]

that is (A,A)-linear. Conversely, the inverse isomorphism \((\mathrm{P}_{1}\otimes_{k}\mathrm{Q}_{1})\otimes_{k}(\mathrm{P}_{2}\otimes_{k}\mathrm{Q}_{2})\to(\mathrm{P}_{1}\otimes_{k}\mathrm{P}_{2})\otimes_{k}(\mathrm{Q}_{1}\otimes_{k}\mathrm{Q}_{2})\) defines an (A,A)-linear morphism \((\mathrm{P}_{1}\otimes_{\mathrm{B}_{1}}\mathrm{Q}_{1})\otimes_{k}(\mathrm{P}_{2}\otimes_{\mathrm{B}_{2}}\mathrm{Q}_{2})\to(\mathrm{P}_{1}\otimes_{k}\mathrm{P}_{2})\otimes_{\mathrm{B}}(\mathrm{Q}_{1}\otimes_{k}\mathrm{Q}_{2})\) . These two morphisms are each other's inverses and are thus isomorphisms. Since the \((\mathrm{A}_{i},\mathrm{A}_{i})\) -bimodule \(P_{i}\otimes_{B_{i}}Q_{i}\) is isomorphic to \(A_{i}\) , we obtain an (A,A)-linear isomorphism \(P\otimes_{B}Q\to A\) . We likewise obtain a (B,B)-linear isomorphism \(Q\otimes_{A}P\to B\) , which completes the proof of d).

Under the assumptions of e), let P be an invertible \((\mathrm{A},\mathrm{B})_{k}\) -bimodule; then \(\mathrm{P}_{(k^{\prime})}\) is an invertible \((\mathrm{A}_{(k^{\prime})},\mathrm{B}_{(k^{\prime})})_{k^{\prime}}\) -bimodule.

Let A be a k-algebra, and let e be an idempotent in A. The set eAe, endowed with the addition, multiplication, and action of k induced by those of A, is a k-algebra with unit element e.

PROPOSITION 14. --- Let A and B be k-algebras. Then A and B are Morita equivalent if and only if there exist an integer \(n \geqslant 1\) and a square matrix \(e = (e_{ij})\) in \(\mathbf{M}_{n}(\mathbf{B})\) satisfying the following conditions:

\begin{enumerate}
\def\labelenumi{(\roman{enumi})}
\item
  We have \(e^2 = e\) .
\item
  The two-sided ideal of B generated by the elements \(e_{ij}\) is equal to B.
\item
  The \(k\) -algebra A is isomorphic to \(e\mathbf{M}_n(\mathrm{B})e\) .
\end{enumerate}

If conditions (i) and (ii) are satisfied, then the \((e\mathbf{M}_n(\mathrm{B})e,\mathrm{B})_k\) -bimodule \(e\mathrm{B}_d^n\) is invertible.

In view of Theorem 1 (VIII, p.~101), the k-algebra A is Morita equivalent to B if and only if it is isomorphic to the endomorphism algebra of a projective, finitely generated, and generating right B-module. The proposition therefore follows from Lemmas 3 and 4 below.

Lemma 3. --- A right B-module P is projective, finitely generated, and generating if and only if there exist an integer \(n \geqslant 0\) and an idempotent \(e = (e_{ij})\) in \(\mathbf{M}_n(\mathrm{B})\) with the following properties:

\begin{enumerate}
\def\labelenumi{(\roman{enumi})}
\item
  The B-module P is isomorphic to \(e\mathrm{B}_d^n\) .
\item
  The two-sided ideal of B generated by the elements \(e_{ij}\) is equal to B.
\end{enumerate}

Let P be a right B-module. Then P is projective and finitely generated if and only if it is isomorphic to a direct factor submodule of a B-module of the form \(B_{d}^{n}\) , where n is an integer \(\geqslant 0\) (II, §2, No.~2, p.~232, Corollary 1). If we identify the k-algebras \(\mathbf{M}_{n}(\mathrm{B})\) and \(\operatorname{End}(\mathrm{B}_{d}^{n})\) , then this means that there exists an idempotent e in \(\mathbf{M}_{n}(\mathrm{B})\) such that P is isomorphic to \(eB_{d}^{n}\) .

The B-module P is generating if and only if its trace ideal \(\tau(\mathrm{P})\) is equal to B, that is, \(\tau(e\mathrm{B}_{d}^{n})=\mathrm{B}\) (VIII, p.~80, Theorem 1). Let \(x_{1},\ldots,x_{n}\) be the elements of \(B_{d}^{n}\) corresponding to the columns of the matrix e, and let \(x_{i}^{*}\) (for \(1\leqslant i\leqslant n\) ) be the linear form \((b_{1},\ldots,b_{n})\mapsto b_{i}\) on \(eB_{d}^{n}\) . The family \((x_{1},\ldots,x_{n})\) generates the B-module \(eB_{d}^{n}\) , and the family \((x_{1}^{*},\ldots,x_{n}^{*})\) generates its dual. Now, we have \(\langle x_{i}^{*},x_{j}\rangle=e_{ij}\) , so \(\tau(eB_{d}^{n})\) is the two-sided ideal of B generated by the \(e_{ij}\) . This proves Lemma 3.

Lemma 4. --- Let V be a B-module and E the endomorphism k-algebra of V. Let e be a projector of V and P the image of e. The mapping that sends \(v \in eEe\) to the endomorphism \(x \mapsto v(x)\) of the B-module P is an isomorphism of k-algebras from eEe to \(\operatorname{End}_{\mathrm{B}}(\mathrm{P})\) .

Denote the mapping described in the statement by \(\varphi:eEe\to\operatorname{End}_{\mathrm{B}}(\mathrm{P})\) ; it is a homomorphism of k-algebras. Let \(u\in\operatorname{End}_{\mathrm{B}}(\mathrm{P})\) . Denote by v the endomorphism of V defined by \(v(x)=u(e(x))\) for \(x\in V\) . We have \((eve)(x)=u(x)\) for \(x\in P\) , that is, \(\varphi(eve)=u\) . Consequently, \(\varphi\) is surjective. Let w be an element of the kernel of \(\varphi\) ; the restrictions of w to the kernel and to the image of e are zero, so w is zero, which proves that \(\varphi\) is injective.

Examples. --- 1) Let A be a k-algebra and e an idempotent in A such that AeA = A. The k-algebra eAe can be identified with the endomorphism k-algebra of the submodule \(eA_{d}\) of \(A_{d}\) (Lemma 4). Since AeA = A, it follows from Proposition 14 that \(eA_{d}\) is an invertible \((eAe, A)_{k}\) -bimodule, so that the \(k\) -algebras eAe and A are Morita equivalent. Moreover, Morita's theorem (VIII, p.~103) implies the following results:

\begin{enumerate}
\def\labelenumi{\alph{enumi})}
\item
  Let M and N be left A-modules. Every eAe-linear mapping from eM to eN extends uniquely to an A-linear mapping from M to N.
\item
  Every left module over the k-algebra eAe is isomorphic to a module of the form eM, where M is a left A-module.
\end{enumerate}

\begin{enumerate}
\def\labelenumi{\arabic{enumi})}
\setcounter{enumi}{1}
\tightlist
\item
  Let A be a k-algebra and \(n \geqslant 1\) an integer. We identify the matrix algebra \(\mathbf{M}_{n}(\mathrm{A})\) with the endomorphism algebra of the right A-module \(A_{d}^{n}\) . We have seen that A and \(\mathbf{M}_{n}(\mathrm{A})\) are Morita equivalent. For every left A-module M, identify \(A_{d}^{n} \otimes_{A} M\) with \(M^{n}\) . The algebra \(\mathbf{M}_{n}(\mathrm{A})\) then has a left action on \(M^{n}\) , and we have
\end{enumerate}

\[
(a \cdot m) _ {i} = \sum_ {j = 1} ^ {n} a _ {i j} m _ {j}
\]

for \(a = (a_{ij})\) in \(\mathbf{M}_n(\mathrm{A})\) and \(m = (m_i)\) in \(\mathrm{M}^n\) . Morita's theorem implies the following results:

\begin{enumerate}
\def\labelenumi{\alph{enumi})}
\item
  Every left module over the algebra \(\mathbf{M}_n(\mathrm{A})\) is isomorphic to a module of the form \(\mathrm{M}^n\) , where \(\mathrm{M}\) is a left A-module.
\item
  Let \(\mathrm{M}\) be a left A-module. The mapping \(\mathrm{N} \mapsto \mathrm{N}^n\) is a bijection from the set of A-submodules of \(\mathrm{M}\) to the set of \(\mathbf{M}_n(\mathrm{A})\) -submodules of \(\mathrm{M}^n\) .
\item
  Let M and N be left A-modules. For an A-linear mapping \(g: \mathrm{M} \to \mathrm{N}\) , let \(g_{n}\) be the mapping \((m_{i}) \mapsto (g(m_{i}))\) from \(\mathrm{M}^{n}\) to \(\mathrm{N}^{n}\) . Then the mapping \(g \mapsto g_{n}\) is a bijection from \(\operatorname{Hom}_{\mathrm{A}}(\mathrm{M}, \mathrm{N})\) to \(\operatorname{Hom}_{\mathbf{M}_{n}(\mathrm{A})}(\mathrm{M}^{n}, \mathrm{N}^{n})\) .
\item
  Let M be a left A-module. The module \(M^{n}\) over the ring \(\mathbf{M}_{n}(\mathrm{A})\) is indecomposable (resp. semisimple, simple, Artinian, Noetherian, finitely generated) if and only if the A-module M is.
\end{enumerate}

\begin{enumerate}
\def\labelenumi{\arabic{enumi})}
\setcounter{enumi}{2}
\tightlist
\item
  Let A be a principal ideal domain and L a nonzero, finitely generated, free A-module. Let B be the endomorphism ring of L; then L is an invertible \((\mathrm{A},\mathrm{B})_{\mathbf{Z}}\) -bimodule, and the rings A and B are Morita equivalent. By Morita's theorem, Proposition 10, a) (VIII, p.~109), and the structure theorem for finitely generated A-modules (VII, §4, No.~4, p.~19, Theorem 2), every finitely generated B-module is isomorphic to \(\oplus_{i=1}^{m}(\mathrm{L}/\mathfrak{a}_{i}\mathrm{L})\) , where m is a natural number and the \(a_{i}\) are ideals of A satisfying \(a_{1} \subset a_{2} \subset \cdots \subset a_{m}\) and \(a_{m} \neq A\) ; the integer m and the ideals \(a_{i}\) are uniquely determined. By Proposition 6 of VIII, p.~105, every two-sided ideal of B is of the form dB, where d is an element of A.
\end{enumerate}

